% META: {"id": "1SPE-SUITES-EV-B", "chapitre": "1SPE-SUITES", "type_objet": "evaluation", "version": "B", "duree_min": 55, "bareme_total": 20, "capacites_codes": ["C1", "C2", "C3", "C4", "C5", "C6", "C7"], "competences": ["chercher", "modeliser", "calculer", "raisonner", "communiquer"], "mode_creation": "ex_nihilo", "sources_inspiration": [], "parametres_sympy": {"u0_arith": 5, "r_arith": 3, "u0_geo": 4, "q_geo": 2, "capital": 1000, "taux_annuel": 4, "seuil": 1500, "p0_theatre": 18, "r_theatre": 4, "nombre_rangees": 12}, "fichier_tex": "chapitres/1SPE-SUITES/evaluations/1SPE-SUITES-EV-B.tex", "corrige_tex": "chapitres/1SPE-SUITES/evaluations/1SPE-SUITES-EV-B-corrige.tex", "status": "generated"}
% BEGIN-VERIFY
% from sympy import Rational, symbols, simplify, log
% n = symbols('n', integer=True, nonnegative=True)
%
% # ---- Exercice 1 : suite arithmetique ----
% u0_arith = 5
% r_arith  = 3
% u_arith  = u0_arith + r_arith * n
% assert u_arith.subs(n, 0) == 5
% assert u_arith.subs(n, 1) == 8
% assert u_arith.subs(n, 2) == 11
% assert u_arith.subs(n, 3) == 14
% # Difference constante
% diff_arith = u_arith.subs(n, n+1) - u_arith
% assert simplify(diff_arith) == 3
% # Variations : r > 0 => suite croissante
% assert r_arith > 0
%
% # ---- Exercice 2 : suite geometrique ----
% u0_geo = 4
% q_geo  = 2
% u_geo  = u0_geo * q_geo**n
% assert u_geo.subs(n, 0) == 4
% assert u_geo.subs(n, 1) == 8
% assert u_geo.subs(n, 2) == 16
% # Quotient constant
% quot = u_geo.subs(n, n+1) / u_geo
% assert simplify(quot - q_geo) == 0
% # Somme des 5 premiers termes (k=0 a 4) : u0*(q^5 - 1)/(q-1)
% S5_geo = sum(u0_geo * q_geo**k for k in range(5))
% assert S5_geo == u0_geo * (q_geo**5 - 1) // (q_geo - 1)
% assert S5_geo == 124
%
% # ---- Exercice 3 : placement a interets composes ----
% capital = 1000      # placement initial en euros
% taux_annuel = 4     # taux annuel en %
% q_place = Rational(104, 100)   # 1 + 4/100
% u_place = capital * q_place**n
% assert u_place.subs(n, 0) == 1000
% assert u_place.subs(n, 1) == 1040
% assert round(float(u_place.subs(n, 2))) == 1082   # 1081.6 arrondi a l'euro
% seuil = 1500
% # Plus petit n tel que u_n >= seuil : 1000 * 1.04^n >= 1500 => 1.04^n >= 1.5
% import math
% n_seuil = math.ceil(math.log(seuil/capital) / math.log(float(q_place)))
% assert n_seuil == 11
% assert float(u_place.subs(n, 10)) < seuil
% assert float(u_place.subs(n, 11)) >= seuil
%
% # ---- Exercice 4 : rangees d'un theatre ----
% p = 18 + 4*n
% assert p.subs(n, 0) == 18
% assert p.subs(n, 1) == 22
% assert simplify(p.subs(n, n+1) - p) == 4
% T = lambda m: (m+1)*(18+2*m)
% # La formule erronee (n+1)*p_n echoue des n=1 : 40 != 44.
% assert T(1) == 40
% assert 2 * p.subs(n, 1) == 44
% for m in range(12):
%     assert sum(18 + 4*k for k in range(m+1)) == T(m)
% # Douze rangees correspondent aux indices 0 a 11.
% assert T(11) == 480
% assert sum(18 + 4*k for k in range(12)) == 480
% END-VERIFY

%---------------------------------------------------------------
% Durée : 55 minutes — Barème : 20 points
% Toute réponse doit être justifiée ; les résultats seuls ne sont pas acceptés.
% La qualité de la rédaction sera prise en compte.
%---------------------------------------------------------------

\begin{evaluation}{1SPE-SUITES-EV-B}{55}

%===============================================================
\section*{Exercice 1 \ifnxVersionProfesseur\hfill (6 points)\fi \hfill Capacités C1, C2, C5}
%===============================================================

\textit{Compétences évaluées : calculer, raisonner, communiquer.}

\medskip
On considère la suite $(u_n)$ définie pour tout entier naturel $n$ par :
\[
  u_0 = 5 \qquad \text{et} \qquad u_{n+1} = u_n + 3.
\]

\begin{enumerate}
  \item Calculer $u_1$, $u_2$ et $u_3$.
  \ifnxVersionProfesseur\hfill\textit{(1 pt — calculer)}\fi

  \item Démontrer que la suite $(u_n)$ est arithmétique. Préciser sa raison $r$ et son premier terme $u_0$.
  \ifnxVersionProfesseur\hfill\textit{(2 pts — raisonner, communiquer)}\fi

  \item Exprimer $u_n$ en fonction de $n$ pour tout entier naturel $n$.
  \ifnxVersionProfesseur\hfill\textit{(1 pt — calculer)}\fi

  \item Étudier les variations de la suite $(u_n)$. Justifier rigoureusement.
  \ifnxVersionProfesseur\hfill\textit{(2 pts — raisonner, communiquer)}\fi
\end{enumerate}

%===============================================================
\section*{Exercice 2 \ifnxVersionProfesseur\hfill (5 points)\fi \hfill Capacités C3, C4}
%===============================================================

\textit{Compétences évaluées : calculer, raisonner, communiquer.}

\medskip
On considère la suite $(u_n)$ définie pour tout entier naturel $n$ par :
\[
  u_n = 4 \times 2^n.
\]

\begin{enumerate}
  \item Calculer $u_0$, $u_1$ et $u_2$.
  \ifnxVersionProfesseur\hfill\textit{(1 pt — calculer)}\fi

  \item Démontrer que la suite $(u_n)$ est géométrique. Préciser sa raison $q$ et son premier terme $u_0$.
  \ifnxVersionProfesseur\hfill\textit{(2 pts — raisonner, communiquer)}\fi

  \item Calculer la somme $S = u_0 + u_1 + u_2 + u_3 + u_4$.

  On rappelle que, pour une suite géométrique de raison $q
\neq 1$, on a :
  \[
    \sum_{k=0}^{n} u_k = u_0 \times \frac{q^{n+1} - 1}{q - 1}.
  \]
  \ifnxVersionProfesseur\hfill\textit{(2 pts — calculer)}\fi
\end{enumerate}

%===============================================================
\section*{Exercice 3 \ifnxVersionProfesseur\hfill (5 points)\fi \hfill Capacités C6, C7}
%===============================================================

\textit{Compétences évaluées : modéliser, chercher, communiquer.}

\medskip
Lucas place $1\,000$\,€ sur un livret d'épargne qui rapporte $4\,\%$ d'intérêts par an. On note $u_n$ le capital disponible (en euros) au bout de $n$ années entières (avec $u_0 = 1\,000$).

\begin{enumerate}
  \item Justifier que, pour tout entier naturel $n$ :
  \[
    u_{n+1} = 1{,}04 \times u_n.
  \]
  \ifnxVersionProfesseur\hfill\textit{(1 pt — modéliser)}\fi

  \item En déduire que $(u_n)$ est une suite géométrique. Exprimer $u_n$ en fonction de $n$.
  \ifnxVersionProfesseur\hfill\textit{(1 pt — calculer)}\fi

  \item Calculer $u_1$ et $u_2$ (valeurs exactes, ou arrondies à l'euro si nécessaire).
  \ifnxVersionProfesseur\hfill\textit{(1 pt — calculer)}\fi

  \item Le programme Python suivant a pour objectif de déterminer le plus petit entier $n$ tel que $u_n \geq 1\,500$.

\lstinputlisting[language=Python]{chapitres/1SPE-SUITES/python/1SPE-SUITES-EV-B-SEUIL.py}

  \begin{enumerate}
    \item Expliquer, ligne par ligne, le rôle de la boucle \texttt{while}.
    \ifnxVersionProfesseur\hfill\textit{(1 pt — communiquer)}\fi
    \item Indiquer la valeur affichée par \texttt{print(n)} et vérifier par le calcul que cette valeur est cohérente.
    \ifnxVersionProfesseur\hfill\textit{(1 pt — chercher)}\fi
  \end{enumerate}
\end{enumerate}

%===============================================================
\section*{Exercice 4 \ifnxVersionProfesseur\hfill (4 points)\fi \hfill Capacités C2, C4}
%===============================================================

\textit{Compétences évaluées : modéliser, calculer, raisonner.}

\medskip
Dans un théâtre, la première rangée compte $18$ places et chaque nouvelle
rangée compte $4$ places de plus que la précédente. On note $p_n$ le nombre de
places de la rangée d'indice $n$ ; ainsi, $n=0$ correspond à la première rangée
et $p_0=18$.

Pour tout entier naturel $n$, on note $T_n$ le nombre total de places des
$(n+1)$ premières rangées :
\[
  T_n=p_0+p_1+\cdots+p_n=\sum_{k=0}^{n}p_k.
\]

\begin{enumerate}
  \item Exprimer $p_{n+1}$ en fonction de $p_n$, puis justifier que $(p_n)$ est
  une suite arithmétique et établir $p_n=18+4n$.
  \ifnxVersionProfesseur\hfill\textit{(1 pt — modéliser, raisonner)}\fi

  \item Lina affirme que $T_n=(n+1)\times p_n$. Réfuter cette formule pour
  $n=1$, puis établir que $T_n=(n+1)(18+2n)$ pour tout entier naturel $n$.
  \ifnxVersionProfesseur\hfill\textit{(2 pts — calculer, raisonner)}\fi

  \item Calculer le nombre total de places des $12$ premières rangées.
  \ifnxVersionProfesseur\hfill\textit{(1 pt — calculer)}\fi
\end{enumerate}

\end{evaluation}
